% ============================================================================
%  Chapter 19 -- Instructions and Interrupts                     (Lecture 19)
% ============================================================================
\chapter{Instructions and Interrupts}
\label{ch:instructions}

\epigraphline{Why does the CPU never get any work done? Because it gets
interrupted constantly.}{Anonymous joke about computer architecture}

\begin{chapgoals}
By the end of this chapter you will be able to
\begin{itemize}
\item describe an instruction as an operation code and its operands, and name the kinds of operand;
\item decode a \riscv instruction by hand, recognising its format and fields;
\item follow an instruction through the processor cycle: fetch, decode, execute, memory, write back;
\item explain what a pipeline is and why it makes a processor faster;
\item define interrupts, exceptions and traps, and list their causes;
\item describe what \riscv does when a trap occurs and how the interrupted program resumes.
\end{itemize}
\end{chapgoals}
% Generated by tools/textbook/export_listings.js from tools/decks/mastery.js and tools/concepts/registry.json: do not edit.
\begin{chapstatus}{Learning status of this chapter}
\masterychip{now}{UNDERSTAND NOW}\enspace instruction formats; addressing modes; the processor cycle; interrupts and traps.\par\smallskip
\masterychip{use}{USE NOW}\enspace the complete program, the PRG32 tools (used, not studied, in this course).\par\smallskip
\masterychip{later}{UNDERSTAND LATER}\enspace nothing.
\end{chapstatus}
\begin{chapstatus}{Roadmap of this chapter}
By the end of the chapter you can:
\begin{enumerate}[nosep,leftmargin=2em]
\item decode a RISC-V instruction into its fields
\item name the addressing mode of each operand
\item describe the instruction cycle and what an interrupt changes
\end{enumerate}
It builds on: only the setup of Chapter~\ref{ch:setup}. Next: Chapter~\ref{ch:video}.
\end{chapstatus}


% ----------------------------------------------------------------------------
\section{Warm-up on the machine}

\txtfile[caption={Four words in memory.},label=lst:ins:words]{code/l19/c03.txt}

Listing~\ref{lst:ins:words} is the function \code{first\_update} of
Chapter~\ref{ch:firstcontact} as it sits in memory: four numbers, four
addresses apart. (The \ins{la} that precedes them is omitted, and the
encodings are the uncompressed ones of 32 bits: with the C extension the
assembler may choose forms of 16 bits for some of them.) Three questions
follow. How does the processor know that the first number is a load and the
second an addition? Which steps does it take for each of them? And while it
runs them, how does it notice that a button was pressed or that a timer has
expired?

% ----------------------------------------------------------------------------
\section{Instructions}

\subsection{What an instruction is made of}

The typical instructions of a processor are transfers of data between memory
and registers, arithmetic and logic operations, instructions for the control
of the flow, and transfers of data to and from input/output. (In \riscv the
fourth family does not exist as such, because devices are memory.)

\begin{definition}{Instruction}
An \term{instruction} is composed of an \term{operation code}, the opcode,
and of zero or more \term{operands}. An operand can be a constant (an
\emph{immediate}), a memory location, or a register; it can be explicit or
implicit.
\end{definition}

Instructions can have zero, one, two or three operands: \ins{add} has three,
all registers; \ins{ret} has none visible.

\subsection{The formats of RISC-V instructions}

Every \riscv instruction of the base set is 32 bits long, and is laid out in
one of six \term{formats} (Figure~\ref{fig:ins:formats}). The formats were
designed so that the same field is always in the same place: the opcode in
bits 6 to 0, the destination register \code{rd} in bits 11 to 7, the source
registers \code{rs1} and \code{rs2} in bits 19 to 15 and 24 to 20, the field
\code{funct3} in bits 14 to 12. The formats B and J differ from S and U only
in how the bits of the immediate are arranged.

\begin{figure}[htb]
\centering\footnotesize
\begin{bytefield}[bitwidth=0.0275\textwidth,bitheight=6.5mm]{32}
\bitheader[endianness=big]{0,6,7,11,12,14,15,19,20,24,25,31}\\
\begin{rightwordgroup}{R}\bitbox{7}[bgcolor=prgviolet!20]{funct7} & \bitbox{5}[bgcolor=prgaccent!20]{rs2} & \bitbox{5}[bgcolor=prgaccent!20]{rs1} & \bitbox{3}[bgcolor=prgviolet!20]{funct3} & \bitbox{5}[bgcolor=prggreen!20]{rd} & \bitbox{7}[bgcolor=prgsteel!20]{opcode}\end{rightwordgroup}\\[2pt]
\begin{rightwordgroup}{I}\bitbox{12}[bgcolor=prgamber!30]{imm[11:0]} & \bitbox{5}[bgcolor=prgaccent!20]{rs1} & \bitbox{3}[bgcolor=prgviolet!20]{funct3} & \bitbox{5}[bgcolor=prggreen!20]{rd} & \bitbox{7}[bgcolor=prgsteel!20]{opcode}\end{rightwordgroup}\\[2pt]
\begin{rightwordgroup}{S}\bitbox{7}[bgcolor=prgamber!30]{imm[11:5]} & \bitbox{5}[bgcolor=prgaccent!20]{rs2} & \bitbox{5}[bgcolor=prgaccent!20]{rs1} & \bitbox{3}[bgcolor=prgviolet!20]{funct3} & \bitbox{5}[bgcolor=prgamber!30]{imm[4:0]} & \bitbox{7}[bgcolor=prgsteel!20]{opcode}\end{rightwordgroup}\\[2pt]
\begin{rightwordgroup}{B}\bitbox{7}[bgcolor=prgamber!30]{imm[12$|$10:5]} & \bitbox{5}[bgcolor=prgaccent!20]{rs2} & \bitbox{5}[bgcolor=prgaccent!20]{rs1} & \bitbox{3}[bgcolor=prgviolet!20]{funct3} & \bitbox{5}[bgcolor=prgamber!30]{imm[4:1$|$11]} & \bitbox{7}[bgcolor=prgsteel!20]{opcode}\end{rightwordgroup}\\[2pt]
\begin{rightwordgroup}{U}\bitbox{20}[bgcolor=prgamber!30]{imm[31:12]} & \bitbox{5}[bgcolor=prggreen!20]{rd} & \bitbox{7}[bgcolor=prgsteel!20]{opcode}\end{rightwordgroup}\\[2pt]
\begin{rightwordgroup}{J}\bitbox{20}[bgcolor=prgamber!30]{imm[20$|$10:1$|$11$|$19:12]} & \bitbox{5}[bgcolor=prggreen!20]{rd} & \bitbox{7}[bgcolor=prgsteel!20]{opcode}\end{rightwordgroup}
\end{bytefield}
\caption{The six formats of \riscv. R: register to register (\ins{add},
\ins{sub}, \ins{and}, \ins{sll}, \ldots). I: immediates and loads
(\ins{addi}, \ins{andi}, \ins{lw}, \ins{jalr}, \ldots). S: stores. B:
branches. U: upper immediates (\ins{lui}, \ins{auipc}). J: jumps
(\ins{jal}).}
\label{fig:ins:formats}
\end{figure}

This regularity is the heart of the RISC idea. The processor always knows
where to look first, and can read the registers named by an instruction while
it is still finding out which instruction it is.

\subsection{Decoding by hand}

\begin{prgconcept}[The method]
Write the eight hexadecimal digits in binary; cut off the last seven bits, the
opcode; look up the format; then cut the fields.
\end{prgconcept}

\begin{example}{\hex{006283B3}}
$\hex{006283B3}=0000\,0000\,0110\,0010\,1000\,0011\,1011\,0011_2$. The last
seven bits, \texttt{0110011}, are the opcode of the register-to-register
operations: the format is R.
\begin{center}\small\ttfamily
\begin{tabular}{|c|c|c|c|c|c|}\hline
0000000&00110&00101&000&00111&0110011\\\hline
\multicolumn{1}{c}{\textrm{\scriptsize funct7}}&\multicolumn{1}{c}{\textrm{\scriptsize rs2 = x6 = t1}}&\multicolumn{1}{c}{\textrm{\scriptsize rs1 = x5 = t0}}&\multicolumn{1}{c}{\textrm{\scriptsize funct3}}&\multicolumn{1}{c}{\textrm{\scriptsize rd = x7 = t2}}&\multicolumn{1}{c}{\textrm{\scriptsize opcode}}
\end{tabular}
\end{center}
With $\code{funct3}=000$ and $\code{funct7}=0000000$ the operation is an
addition (with $\code{funct7}=0100000$ it would be a subtraction): the
instruction is \code{add t2, t0, t1}.
\end{example}

\begin{example}{\hex{0002A303}}
$\hex{0002A303}=0000\,0000\,0000\,0010\,1010\,0011\,0000\,0011_2$. The opcode
\texttt{0000011} is that of the loads: the format is I.
\begin{center}\small\ttfamily
\begin{tabular}{|c|c|c|c|c|}\hline
000000000000&00101&010&00110&0000011\\\hline
\multicolumn{1}{c}{\textrm{\scriptsize immediate = 0}}&\multicolumn{1}{c}{\textrm{\scriptsize rs1 = x5 = t0}}&\multicolumn{1}{c}{\textrm{\scriptsize funct3 = word}}&\multicolumn{1}{c}{\textrm{\scriptsize rd = x6 = t1}}&\multicolumn{1}{c}{\textrm{\scriptsize opcode}}
\end{tabular}
\end{center}
The instruction is \code{lw t1, 0(t0)}, the first of the warm-up.
\end{example}

The decoders of Chapter~\ref{ch:combinational} do the same thing in a fraction
of a nanosecond. The tables needed to decode any instruction are in
Appendix~\ref{app:riscv}.

\subsection{Instructions of fixed or of variable length}

In a machine with instructions of \emph{fixed length} all the instructions
have the same number of bits, the next instruction is always at the PC plus a
constant, and fetch and decode are simple. RISC processors are of this kind:
\riscv uses 32 bits and, with the compressed extension, 16 or 32, told apart
by the two lowest bits (Chapter~\ref{ch:combinational}). In a machine with
instructions of \emph{variable length} the length is known only after the
first part has been decoded: programs are compact, and decoding is complex.
CISC processors are of this kind: the MC68000 uses one to five words of 16
bits, and Intel x86 one to fifteen bytes. In either case the length of the
instruction decides by how much the PC advances.

\subsection{Addressing modes: where an operand is found}

\begin{table}[htb]
\centering\small
\caption{Addressing modes.}\label{tab:ins:modes}
\begin{tabularx}{\textwidth}{@{}lLl@{}}
\toprule
\textbf{Mode} & \textbf{The operand is} & \textbf{\riscv example}\\
\midrule
immediate & a constant inside the instruction & \code{addi t1, t1, 1}\\
register & the content of a register & \code{add t2, t0, t1}\\
base plus displacement & the memory word at register + constant & \code{lw t1, 8(sp)}\\
PC-relative & an address at a distance from the instruction itself & \code{beq t0, t1, loop}\\
memory indirect & the memory word whose address is in memory & not available: two loads\\
\bottomrule
\end{tabularx}
\end{table}

\riscv has only the first four modes of Table~\ref{tab:ins:modes}. A CISC
processor such as the MC68000 has a dozen, including memory-indirect and
auto-increment modes. The operation $R:=M[M[A]]$ becomes in \riscv
\code{lw t0, 0(t1)} followed by \code{lw t0, 0(t0)}. The price of rich
addressing modes is a further phase in the processor cycle, in which the
operands are assembled.

\subsection{Calls, returns and the calling convention}

The instructions for the control of the flow deserve one more word, because
every function of \vermone depends on them. \ins{jal} (\emph{jump and link},
format J) jumps to an address relative to the PC and writes in its destination
register the address of the following instruction; \ins{jalr} (format I)
does the same with a target taken from a register. The pseudo-instruction
\ins{call} is a \ins{jal} (or an \ins{auipc} with a \ins{jalr}, for distant
targets) whose destination is \reg{ra}; \ins{ret} is \code{jalr zero, 0(ra)}.

The \term{calling convention} is the agreement that lets separately written
functions call one another \citep{riscvabi}. It says that arguments travel in
\reg{a0}--\reg{a7} and results in \reg{a0} and \reg{a1}; that a function may
destroy the registers \reg{t0}--\reg{t6} and \reg{a0}--\reg{a7}, which are
therefore \emph{caller-saved}; that it must return the registers
\reg{s0}--\reg{s11} and \reg{sp} as it found them, which are therefore
\emph{callee-saved}; and that \reg{sp} is always a multiple of 16. This is
why, since Chapter~\ref{ch:verm1p2}, a value that must survive a call has been
kept in an \reg{s} register and saved in the frame. The convention is stated
in full in Appendix~\ref{app:riscv}.

% ----------------------------------------------------------------------------
\section{The processor cycle}

\subsection{The phases of the cycle}

After a start-up phase, the \emph{bootstrap}, a processor repeats for ever the
same sequence: its fundamental algorithm. In the classic RISC formulation
the sequence has five phases (Figure~\ref{fig:ins:cycle}).

\begin{figure}[htb]
\centering
\begin{tikzpicture}[node distance=4mm]
  \node[blkacc,text width=21mm,minimum height=19mm] (a) {\textbf{Fetch}\\[1pt]\footnotesize read the instruction; the PC advances};
  \node[blkamb,text width=21mm,minimum height=19mm,right=of a] (b) {\textbf{Decode}\\[1pt]\footnotesize recognise it; read the registers};
  \node[blkvio,text width=21mm,minimum height=19mm,right=of b] (c) {\textbf{Execute}\\[1pt]\footnotesize the ALU computes};
  \node[blkgrn,text width=21mm,minimum height=19mm,right=of c] (d) {\textbf{Memory}\\[1pt]\footnotesize loads read, stores write};
  \node[blkred,text width=21mm,minimum height=19mm,right=of d] (e) {\textbf{Write back}\\[1pt]\footnotesize the result goes to \texttt{rd}};
  \foreach \x/\y in {a/b,b/c,c/d,d/e} \draw[->,thick] (\x) -- (\y);
  \draw[->,thick,rounded corners=3mm] (e.south) -- ++(0,-6mm) -| (a.south);
\end{tikzpicture}
\caption{The processor cycle.}
\label{fig:ins:cycle}
\end{figure}

\begin{description}[style=nextline,leftmargin=1.5em]
\item[Fetch.] The instruction at the address in the PC is read from memory;
  the PC advances.
\item[Decode.] The opcode is recognised; the registers named by the
  instruction are read.
\item[Execute.] The ALU computes: a result, an address or a comparison.
\item[Memory.] Loads read the memory and stores write it; the other
  instructions skip this phase.
\item[Write back.] The result is written into the destination register.
\end{description}

\subsection{The fetch phase, transfer by transfer}

In terms of the reference model of Chapter~\ref{ch:processor}, the fetch is
four \term{register transfers} on the internal bus
(Table~\ref{tab:ins:fetch}). In RV32 a word is four bytes and addresses
count bytes, so that the PC advances by 4 (by 2 after a compressed
instruction).

\begin{table}[htb]
\centering\small
\caption{The fetch phase.}\label{tab:ins:fetch}
\begin{tabularx}{\textwidth}{@{}clL@{}}
\toprule
\textbf{Step} & \textbf{Register transfer} & \textbf{Meaning}\\
\midrule
1 & $MA := PC$ & the address of the instruction goes to the memory address register\\
2 & read enable; $MB := M[MA]$ & the memory puts the instruction on the data bus; it is latched in the buffer\\
3 & $D := MB$ & the instruction is copied into the decoding register\\
4 & $PC := PC+4$ & the PC now points to the next instruction\\
\bottomrule
\end{tabularx}
\end{table}

\subsection{The rest of the cycle for two instructions}

\begin{table}[htb]
\centering\small
\caption{The remaining phases for an addition and for a load.}\label{tab:ins:rest}
\begin{tabular}{@{}lll@{}}
\toprule
\textbf{Phase} & \textbf{\code{add t2, t0, t1}} & \textbf{\code{lw t1, 0(t0)}}\\
\midrule
decode & $O_1:=R[\reg{t0}]$;\quad $O_2:=R[\reg{t1}]$ & $O_1:=R[\reg{t0}]$;\quad $O_2:=0$ (the immediate)\\
execute & $U:=O_1+O_2$ & $U:=O_1+O_2$ (an address)\\
memory & -- & $MA:=U$; read enable; $MB:=M[MA]$\\
write back & $R[\reg{t2}]:=U$ & $R[\reg{t1}]:=MB$\\
\bottomrule
\end{tabular}
\end{table}

The same ALU adds two registers for \ins{add} and forms an address for
\ins{lw} (Table~\ref{tab:ins:rest}): only the control signals differ. The
control unit, by decoding the opcode, decides which transfers happen. A store
ends at the memory phase; a branch ends at the execute phase, by choosing the
next PC.

\subsection{The pipeline}

Five phases are five pieces of hardware. While one instruction is executed,
the next can be decoded and a third fetched: this is a \term{pipeline}
(Table~\ref{tab:ins:pipe}).

\begin{table}[htb]
\centering\small
\caption{A linear pipeline of five stages: IF instruction fetch, ID decode,
EX execute, MEM memory, WB write back.}\label{tab:ins:pipe}
\begin{tabular}{@{}l*{8}{>{\centering\arraybackslash}p{8mm}}@{}}
\toprule
\textbf{clock cycle} & 1 & 2 & 3 & 4 & 5 & 6 & 7 & 8\\
\midrule
instruction 1 & \cellcolor{prgaccent!20}IF & \cellcolor{prgamber!25}ID & \cellcolor{prgviolet!18}EX & \cellcolor{prggreen!20}MEM & \cellcolor{prgred!15}WB &&&\\
instruction 2 && \cellcolor{prgaccent!20}IF & \cellcolor{prgamber!25}ID & \cellcolor{prgviolet!18}EX & \cellcolor{prggreen!20}MEM & \cellcolor{prgred!15}WB &&\\
instruction 3 &&& \cellcolor{prgaccent!20}IF & \cellcolor{prgamber!25}ID & \cellcolor{prgviolet!18}EX & \cellcolor{prggreen!20}MEM & \cellcolor{prgred!15}WB &\\
instruction 4 &&&& \cellcolor{prgaccent!20}IF & \cellcolor{prgamber!25}ID & \cellcolor{prgviolet!18}EX & \cellcolor{prggreen!20}MEM & \cellcolor{prgred!15}WB\\
\bottomrule
\end{tabular}
\end{table}

Each instruction still takes five cycles, but one instruction completes at
every clock cycle: the throughput is five times as large. In general, $n$
instructions on a pipeline of $k$ stages complete in $k+n-1$ cycles.
\term{Superscalar} processors put two or more pipelines side by side, and
complete more than one instruction per cycle (Chapter~\ref{ch:beyond}). The
main core of the ESP32-C6 has a pipeline of four stages. Branches, and
dependences between neighbouring instructions, complicate the picture: they
are called \emph{hazards}, and are the subject of an advanced course.

% ----------------------------------------------------------------------------
\section{Interrupts}

\subsection{The problem with the simple cycle}

If the processor only ever did fetch, decode and execute, three problems would
arise. The interaction with the outside world through the I/O units would be
inefficient, because the program would have to keep asking. There would be no
way of forcibly removing an application that, by mistake, enters an infinite
loop: who takes the processor away from a program stuck in \code{j loop}? And
the operating system would have limited control over the system.

\subsection{Interrupts, exceptions, traps}

\begin{definition}{Interrupt}
\term{Interrupts} are exceptional events, usually asynchronous with respect to
the execution of the program in progress, that require the attention of the
processor.
\end{definition}

The \riscv specification uses three words precisely \citep{riscvpriv}. An
\term{interrupt} comes from outside the program, at any moment: a timer, a
device. An \term{exception} is caused by the instruction being executed: an
illegal opcode, a forbidden address. A \term{trap} is the transfer of control
to a handler that follows either of them. (Older processors call ``traps''
the synchronous events and use ``interrupt'' for everything.)

\subsection{The real processor cycle}

The real cycle therefore has one more step (Figure~\ref{fig:ins:real}): after
each instruction the processor checks whether an interrupt is pending and, if
so, passes to the handler instead of fetching the next instruction.

\begin{figure}[htb]
\centering
\begin{tikzpicture}
  \node[blk,fill=prgsteel,text=white,minimum width=22mm,minimum height=9mm] (b) at (0,0) {bootstrap};
  \node[blkacc,minimum width=22mm,minimum height=9mm] (f) at (3.4,0) {fetch};
  \node[blkamb,minimum width=22mm,minimum height=9mm] (x) at (6.8,0) {decode and\\execute};
  \node[blkred,minimum width=22mm,minimum height=10mm] (q) at (6.8,-1.9) {interrupt\\pending?};
  \node[blkvio,minimum width=22mm,minimum height=10mm] (t) at (6.8,-3.8) {trap: go to\\the handler};
  \draw[->,thick] (b) -- (f); \draw[->,thick] (f) -- (x); \draw[->,thick] (x) -- (q);
  \draw[->,thick] (q.west) -| (f.south) node[pos=0.25,above,lbl]{no};
  \draw[->,thick,prgred] (q) -- (t) node[midway,right,lbl,text=prgred]{yes};
  \draw[->,thick,prgviolet] (t.west) -| (f.south);
  \node[lbl,text=prgviolet] at (4.4,-4.1) {the next fetch is at the address in \texttt{mtvec}};
\end{tikzpicture}
\caption{The real processor cycle.}
\label{fig:ins:real}
\end{figure}

\subsection{Possible causes of interruption}

\begin{description}[style=nextline,leftmargin=1.5em]
\item[Periodic interrupt.] A timer, for example every 10~ms: to share the
  processor among programs and to measure the time that each one uses.
\item[I/O interrupt.] A peripheral reports its state, for example that it is
  ready to receive data, to synchronise with the processor.
\item[Errors in the program.] An instruction that does not exist or is
  privileged, a forbidden address: synchronous events, the exceptions. (In
  \riscv an integer overflow does \emph{not} raise an exception.)
\item[Faults.] Failures of the system detected by dedicated probes: power,
  temperature, memory errors.
\item[Reset.] Brings the whole system back to a known state.
\item[Software interrupt.] Generated by a program that needs a shared resource
  and asks the operating system to mediate: in \riscv, the instruction
  \ins{ecall}, \emph{environment call}, which is synchronous with the program
  that executes it.
\end{description}

\subsection{Handling an interrupt}

\begin{enumerate}
\item Save the context, hardware part: at least the address at which to resume.
\item Identify the cause of the interrupt.
\item Jump to the entry point of the \term{interrupt service routine} (ISR).
\item Save the context, software part: the registers that the routine is going
  to use.
\item Serve the interrupt.
\item If the interrupt allows it, restore the context and resume the suspended
  program from the instruction that follows the last one executed.
\end{enumerate}

The \term{context} to save and restore is the PC, the status and the
registers. The saving can be done partly or wholly in hardware or in
software; \riscv does the minimum in hardware.

\subsection{Who interrupted? Two solutions}

With several devices the processor must identify who interrupted. With
\emph{status registers and polling}, each device has a status register: a
device that requests an interrupt sets a bit there and raises a common request
line; the service routine interrogates the devices in order of priority and
serves the first that it finds with its bit high. The scheme is simple to
build but takes time, because it also interrogates devices that asked for
nothing. With \term{vectored interrupts} the identity of the source is
supplied to the processor at the moment of the request, and the processor
jumps directly to the routine that corresponds to it: a very efficient
mechanism, in which the number of bits of the vector limits the number of
different sources.

The \term{interrupt latency} is the delay between the instant in which the
request line goes high and the instant in which the service routine starts
doing useful work; it includes the end of the current instruction, the check
and the switch.

% ----------------------------------------------------------------------------
\section{The RISC-V solution}

\riscv solves the whole problem with a few registers and one instruction.

\subsection{The registers of the trap mechanism}

\begin{table}[htb]
\centering\small
\caption{The machine-mode CSRs of the trap mechanism.}\label{tab:ins:csr}
\begin{tabularx}{\textwidth}{@{}lP{0.3\textwidth}L@{}}
\toprule
\textbf{CSR} & \textbf{Written by} & \textbf{Content}\\
\midrule
\reg{mtvec} & the firmware, once, at start-up & the address of the trap handler, and whether the vectored mode is used\\
\reg{mstatus} & hardware and firmware & MIE, the global interrupt enable, and MPIE, its previous value\\
\reg{mie} & the firmware & one enable bit for each source: software, timer, external\\
\reg{mip} & the hardware & one pending bit for each source\\
\reg{mepc} & the hardware, at each trap & the address of the interrupted instruction\\
\reg{mcause} & the hardware, at each trap & one bit that says interrupt or exception, and a code for the cause\\
\bottomrule
\end{tabularx}
\end{table}

An interrupt is taken only if its bit is set in both \reg{mip} and \reg{mie},
and MIE is 1 (Table~\ref{tab:ins:csr}).

\subsection{What the hardware does on a trap}

\begin{enumerate}
\item $\reg{mepc} :=$ the address of the interrupted instruction.
\item $\reg{mcause} :=$ the cause: the top bit is 1 for an interrupt and 0 for
  an exception, and the other bits hold a code.
\item In \reg{mstatus}: $\mathrm{MPIE}:=\mathrm{MIE}$, then
  $\mathrm{MIE}:=0$, so that further interrupts are disabled.
\item The processor passes to machine mode.
\item $PC :=$ the address in \reg{mtvec}; or, in vectored mode and for
  interrupts, that address plus 4 times the cause.
\end{enumerate}

At the end of the handler the instruction \ins{mret} undoes it all:
$\mathrm{MIE}:=\mathrm{MPIE}$ and $PC:=\reg{mepc}$.

Nothing is pushed on a stack by the hardware: saving the general registers is
left to the handler. The ``hardware context save'' of the general scheme is
reduced to two registers and one bit.

\subsection{The cause of a trap}

\begin{table}[htb]
\centering\small
\caption{Some standard values of \reg{mcause}; chips add their own.}\label{tab:ins:cause}
\begin{tabularx}{\textwidth}{@{}lcL@{}}
\toprule
\textbf{Kind} & \textbf{Code} & \textbf{Cause}\\
\midrule
interrupt & 3 & machine software interrupt\\
interrupt & 7 & machine timer interrupt\\
interrupt & 11 & machine external interrupt: a device, through the interrupt controller\\
exception & 2 & illegal instruction\\
exception & 3 & breakpoint (\ins{ebreak}): how a debugger stops a program\\
exception & 5, 7 & load or store access fault: a forbidden address\\
exception & 8, 11 & environment call (\ins{ecall}) from user mode, from machine mode\\
\bottomrule
\end{tabularx}
\end{table}

The same code 3 means different things for interrupts and for exceptions
(Table~\ref{tab:ins:cause}): the top bit of \reg{mcause} distinguishes them.
With \reg{mtvec} in vectored mode, a timer interrupt enters at
$\mathit{base}+4\times 7$.

\subsection{A trap handler, in outline}

\asmfile[caption={The outline of a trap handler.},label=lst:ins:handler]{code/l19/c07.S}

The interrupted program must not notice anything: every register that the
handler touches is saved and restored (Listing~\ref{lst:ins:handler}). Read
as a signed number, \reg{mcause} is negative for interrupts, so that
\ins{bgez} separates the two families. \ins{mret} is the only way back: it
re-enables the interrupts and reloads the PC in one step. The instruction
\code{csrr t0, mcause} is the pseudo-instruction for
\code{csrrs t0, mcause, zero}.

The listing is an outline, not the source of \prg: a real handler saves all
the caller-saved registers, usually on a separate stack reached through the
register \reg{mscratch}. The firmware of the kit contains exactly such a
routine, and a cartridge is interrupted many times per frame.

\subsection{Priorities and nesting}

To handle \term{nested} interrupts a priority is established: while a service
routine runs, only requests from devices of higher priority are accepted, and
a request from a device of high priority, such as the clock, must be accepted
even while one of lower priority is being served. Table~\ref{tab:ins:prio}
compares the classical solution of the MC68000 with that of \riscv, where
priorities belong to an interrupt controller outside the core (on Espressif
chips they go from 1 to 15).

\begin{table}[htb]
\centering\small
\caption{Priorities and nesting in two processors.}\label{tab:ins:prio}
\begin{tabularx}{\textwidth}{@{}P{0.22\textwidth}LL@{}}
\toprule
 & \textbf{Motorola MC68000} & \textbf{\riscv}\\
\midrule
priority of a request & three lines encode an interrupt priority level, 1 to 7 & the interrupt controller gives each source a priority\\
priority of the processor & three bits in the status register: the mask & a threshold register in the interrupt controller\\
a request is accepted if & its level is higher than the mask & its priority is higher than the threshold, and it is enabled\\
while serving it & the mask is raised to that level & interrupts are disabled ($\mathrm{MIE}=0$) unless the handler re-enables them\\
cannot be masked & level 7 & a non-maskable interrupt, defined by each chip\\
identity of the source & an 8-bit vector number from the device: 256 vectors & \reg{mcause}, then the controller says which device\\
return & \code{RTE}, which pops PC and status from the stack & \ins{mret}, which uses \reg{mepc} and \reg{mstatus}\\
\bottomrule
\end{tabularx}
\end{table}

\begin{prgwarn}
It is sometimes said that ``an interrupt is a signal that cannot be masked''.
The opposite is true: most interrupts can be masked, and only a few cannot.
\end{prgwarn}

\subsection{Interrupts on the PRG32 kit}

The three problems of the simple cycle, and their solution, are all visible
on the kit.

\begin{itemize}
\item A timer interrupts the processor a hundred times per second: the
  firmware uses it to keep time and to pace the frames.
\item The display and sound peripherals interrupt when a block of data has
  been sent and the next is wanted.
\item The Wi-Fi radio interrupts when a packet arrives: a cartridge can be
  uploaded while a game is running.
\item The functions \code{init}, \code{update} and \code{draw} of a cartridge
  are ordinary functions called by the firmware between one interrupt and the
  next.
\item A cartridge stuck in a loop does not stop the interrupts. \prgqt ends
  the frame with an error (``instruction budget exhausted''); on the board a
  watchdog reports the stall, and the user presses reset.
\end{itemize}

In \prgqt a breakpoint is the exception with code 3.

% ----------------------------------------------------------------------------
% Generated by tools/textbook/export_listings.js from decks/plan.js and code/l19.S: do not edit.
\section{The program of this chapter}\label{sec:lab:l19}

The listings of this chapter are parts of one program, \fname{code/l19.S} in the
repository of the course, printed in full in Listing~\ref{lst:lab:l19}. The
automatic tests of the course material build it and run it on the headless
runner of \prgqt{}, with scripted input, and check the frame;
Figure~\ref{fig:lab:l19} is the frame that the test wrote.

\begin{description}[style=nextline,leftmargin=1.5em]
\item[On the screen.] Four rows of 32 lamps, the machine code of lw t1, 0(t0) / addi t1, t1, 1 / sw t1, 0(t0) / ret, and beside each a green square: the word in memory is the one decoded by hand in the lecture.
\item[What to press.] Nothing: the program runs by itself.
\end{description}

\begin{figure}[htb]
\centering
\includegraphics[width=0.5\textwidth]{lab_l19.png}
\caption{\fname{l19.S} on \prgqt{}.}\label{fig:lab:l19}
\end{figure}

\asmfile[caption={The complete program of the chapter, \fname{code/l19.S}.},label=lst:lab:l19]{code/l19/program.S}

To run it on \prgqt{} the commands are those of Listing~\ref{lst:lab:l19:run};
Section~\ref{sec:prg:testing} gives the same steps for QEMU and for the board.

\begin{shell}[caption={Building \fname{l19.S} and running it on \prgqt{}.},label=lst:lab:l19:run]
# 1. start PRG32-QT on this computer (it listens on port 8080)

# 2. in the PRG32 folder, with ESP-IDF active: build the cartridge
python3 -m prg32 cartridge build ~/arc12/code/l19.S --portable \
    --entry-prefix l19 --name l19 --out l19.prg32

# 3. send it to PRG32-QT and start it
python3 -m prg32 esp32c6 upload-and-run l19.prg32 \
    --url http://localhost:8080 --slot cart0

# 4. look, press, compare with the previous slide
\end{shell}


% ----------------------------------------------------------------------------
\section{Check your understanding}

\begin{quiz}{In which bits of a \riscv instruction is the opcode?}
\item the seven most significant
\item the seven least significant
\item bits 14 to 12
\item it depends on the format
\end{quiz}

\begin{quiz}{During the fetch of an RV32 instruction that is not compressed, the PC is increased by:}
\item 1
\item 2
\item 4
\item 32
\end{quiz}

\begin{quiz}{Which register holds the address at which an interrupted \riscv program will resume?}
\item \reg{ra}
\item \reg{mtvec}
\item \reg{mepc}
\item \reg{mcause}
\end{quiz}

\begin{quiz}{Which of these is an exception rather than an interrupt, in \riscv terms?}
\item a timer expires
\item a key is pressed
\item an illegal instruction is fetched
\item a packet arrives from the network
\end{quiz}

\section{Exercises}

\begin{exercise}[paper first, then \prg]
\begin{enumerate}[nosep]
\item Decode \hex{00A28293} and \hex{40628333} by hand: format, fields,
  assembly.
\item Encode \code{sub t0, t1, t2} and \code{sw a0, 8(sp)}.
\item Write the register transfers of the whole cycle for
  \code{sw t1, 0(t0)} and for \code{beq t0, t1, target}.
\item Draw the pipeline diagram for six instructions on a pipeline of four
  stages: when does the sixth complete?
\item In \prgqt set a breakpoint in \code{first\_update}, step, and compare
  the machine code shown with your decoding of the warm-up.
\end{enumerate}
\expected{On paper: \code{addi t0, t0, 10} (immediate 10, \code{rs1} =
\reg{x5}, \code{funct3} = 000, \code{rd} = \reg{x5}, opcode 0010011) and
\code{sub t1, t0, t1} (\code{funct7} = 0100000, \code{rs2} = \reg{x6},
\code{rs1} = \reg{x5}, \code{rd} = \reg{x6}); \hex{407302B3} and
\hex{00A12423}; the sixth instruction completes at cycle 9. In \prgqt: the
four words of the warm-up, or their compressed equivalents.}
\end{exercise}

\section{Glossary}

\begin{glossarytable}
operation code (opcode) & codice operativo\\
operand & operando\\
addressing mode & modo di indirizzamento\\
fetch / decode / execute & prelievo / decodifica / esecuzione\\
pipeline stage & stadio della pipeline\\
context & contesto\\
interrupt & interruzione\\
exception / trap & eccezione / trap\\
interrupt service routine & routine di servizio dell'interruzione\\
interrupt latency & latenza di interruzione\\
to mask an interrupt & mascherare un'interruzione\\
nesting & annidamento\\
\end{glossarytable}

\section{Conclusions and reading}

An instruction is an operation code and zero, one or more operands; in \riscv
it is always 32 bits, in six regular formats. The length of the instructions,
fixed or variable, determines how the PC behaves. The processor has a cycle:
it fetches an instruction, decodes it, executes it and writes the result.
Exceptional events break the cycle: the processor saves where it was, serves
the event and resumes. Interrupts can be enabled, masked and ordered by
priority; \riscv does it with a handful of registers and the instruction
\ins{mret}.

\paragraph{Reading.} \citet{patterson2020}, sections 2.5, 4.1--4.6 and 4.10;
the \riscv instruction set manual, volume~II, ``Privileged Architecture''
\citep{riscvpriv}; \emph{Playing with RISC-V}, chapter~4, on \prgqt and its
debugger \citep{prg32textbook}.
